% A2-families.tex -- fragment of appendix-synthobj.tex.
% Heading and label are fixed; a later task adds subsections beneath and
% never renames the heading or the label.
\SOlevelB{The fourteen study variants}\label{app:so:a2}

The construction in~\ref{app:so:a1} builds one component; this section names
the $\SOnumVariants$ assembled objectives the thesis actually measures on. Each is
one entry of the \SOcode{FAMILIES} registry in \SOfile{synthobj/families.py} ---
a frozen \SOcode{FamilySpec} giving the number of active coordinates, their
shares and lengthscales, and whether they carry an interaction, a rotation, a
non-default generator or the monotone rejection wrapper --- and
\SOcode{STUDY\_GRID} lists them in the order the entries below follow. Ten
seeds per variant give the $\SOnumFiles$-file archived grid
(\ref{app:so:a3:files}); every regret curve in the thesis is measured on one
row of Table~\ref{app:so:a2:tab:families}.

A variant fixes shape, never position. \SOcode{select\_active} takes a fresh
permutation of the $D$ coordinates from the seed's own selection stream and
keeps the first $|S|$ of them, so which coordinates are active is redrawn per
seed and index position cannot leak into any result. Because that permutation
is a pure function of the stream and does not depend on $|S|$, truncating it to
different lengths at one seed \emph{nests}: \SOcode{aligned3}'s coordinates are
the first three of \SOcode{aligned10}'s at the same seed. Pairs are drawn the
same way one level down, by \SOcode{select\_pairs}, which cuts disjoint pairs of
\emph{ranks} from a permutation of the active ranks and leaves the caller to map
them through $S$; they are drawn identically whether they go on to become
interactions or rotations. The full seeding scheme is
in~\ref{app:so:a3:seeding}.

Two places in the thesis's design brief have been overtaken by the
implementation, and in both the code is authoritative.
Table~\ref{app:so:a2:tab:families} is generated from \SOcode{FAMILIES} and
\SOcode{STUDY\_GRID} themselves, so what it prints is what the study ran; the
one variant whose parameters the brief's own family table no longer matches
says so in its entry. The second is a definition rather than a parameter: the
brief defines a coordinate as active when its share exceeds
$\activeeps=\SOnumActiveEps$, but \SOcode{Labels.active} in the package is set
membership, $s_i>0$, and $\activeeps$ is recorded alongside it in
\SOcode{Labels.active\_eps} for downstream scoring rather than enforced at build
time. The brief's rule is the scoring definition; the label is the ground truth
about which coordinates the objective actually varies in. Only
\SOcode{dense\_weak} makes the two disagree, and that disagreement is its point.

Each entry below has three paragraphs: the parameters, read off
Table~\ref{app:so:a2:tab:families}; what the variant isolates, in the design
brief's own terms; and what the archived grid measured on it.

\begin{table}[htbp]
  \centering
  \caption[The study variants and their parameters]{
  The $\SOnumVariants$ study variants and their parameters. Generated from
  \SOcode{FAMILIES} and \SOcode{STUDY\_GRID} by \SOcode{regen\_tables.py}, so
  no parameter printed here can drift from the specification the study built.
  Exact: every column is a field of the \SOcode{FamilySpec} the registry
  returns, or --- for $\knots$ --- the length of the grid that spec's components
  are drawn on; no draw enters. Protocol, from the file's own header comment:
  built at $D=\SOnumTableD$ per variant, ``family shape is $D$-invariant;
  \SOcode{dense\_weak}'s $|S|$ equals $D$ here''. Every row but
  \SOcode{dense\_weak}'s is therefore the study's configuration verbatim at its
  own $D=\SOnumStudyD$; \SOcode{dense\_weak}'s $|S|$ and its share scale with
  $D$, as its entry explains.}\label{app:so:a2:tab:families}
  \tiny\setlength{\tabcolsep}{1pt}
  \input{\SOroot/tables/tab_families}
\end{table}

\begin{table}[htbp]
  \centering
  \caption[The archived grid, variant by variant]{
  What the archived grid measures on each variant: $\knots$, the spread of
  $\fstar$ over the seeds, and the two assembly invariants ---
  $\Varnu f=1$ and $\Enu f=0$ --- checked per objective. Draw-dependent:
  $\fstar$ and both moment columns move with the draw, so the numbers are those
  of this archive and not of a re-run elsewhere. Measured on the archived grid,
  seeds $0$--$9$, $D=\SOnumStudyD$, this machine's BLAS\@. Protocol, from the
  file's own header comment: the $\Varnu$ and $\Enu$ columns are QMC estimates
  of the \emph{true} $\refnu$-moments --- a fresh scrambled Sobol' sample per
  objective --- ``not the $64$-node $\nuhat$ quadrature ($\Varhat$/$\Ehat$)
  \SOfile{tab\_orthogonality} uses''. Across the whole grid the worst deviation
  of $\Varnu f$ from one is $\SOnumWorstVar$ and the worst $|\Enu f|$ is
  $\SOnumWorstMean$. The final row repeats one rotated objective at a
  sixteen-fold larger Sobol' sample, so the residual visible in the
  \SOcode{rotated\_t45} row above can be read as sampling error rather than a
  defect in the objective.}\label{app:so:a2:tab:grid}
  \footnotesize\setlength{\tabcolsep}{3pt}
  \input{\SOroot/tables/tab_grid_summary}
\end{table}

\SOlevelC{\SOcode{aligned3}}\label{app:so:a2:aligned3}

\textbf{Parameters.} Three active coordinates carrying equal shares $1/|S|$,
one lengthscale for all of them, no pairs, $\gamma=0$, no rotation, the
Mat\'ern generator, no monotone rejection, and the unit-interval grid
($\knots=\SOnumKnotsUnit$); Table~\ref{app:so:a2:tab:families} gives the
values. $S$ is a fresh random subset at each seed, and its three coordinates
are the first three \SOcode{aligned10} activates at the same seed, in the rank
order shares and lengthscales are assigned in --- the nesting
of~\ref{app:so:a3:seeding}. So the pair \SOcode{aligned3}/\SOcode{aligned10} is
a within-seed comparison of two sparsity levels over a shared coordinate set,
not two unrelated problems.

\textbf{What it isolates.} The design brief's ``positive control: both columns
should agree''. With one lengthscale shared by every active coordinate the
slope factor $5/(3\ell_i^2\vell{\ell_i})$
of~\eqref{app:so:a1:eq:slopeshare} is the same for all $i$ and cancels in the
normalization, so the theoretical slope share is \emph{equal} to the variance
share, coordinate by coordinate --- the two senses of relevance the thesis
contrasts do not merely correlate here, they coincide. A prior whose two
columns disagree on this variant is disagreeing about nothing, so it is the
variant that says whether a difference seen elsewhere is a real effect of the
ground truth or an artefact of the readout.

\textbf{Measured on the archived grid.} Its row of
Table~\ref{app:so:a2:tab:grid} gives the spread of $\fstar$ across the seeds
and the two invariants. Having the smallest active set in the study, it has the
lowest $\fstar$ median in the table.

\SOlevelC{\SOcode{aligned10}}\label{app:so:a2:aligned10}

\textbf{Parameters.} As \SOcode{aligned3} in every field of the specification
except $|S|$: ten active coordinates, equal shares $1/|S|$, the same single
lengthscale, no pairs, $\gamma=0$, Mat\'ern, not monotone, unit grid. Its
active set at a seed contains \SOcode{aligned3}'s.

\textbf{What it isolates.} The same positive control at the sparsity level the
rest of the study uses --- ten active coordinates out of
$D=\SOnumStudyD$ --- so it is the baseline every other ten-coordinate variant
below is read against. Because the shares are equal, no coordinate is easier to
find than another, and a model's recovered support can be scored as a set
without weighting.

\textbf{Measured on the archived grid.} Its row of
Table~\ref{app:so:a2:tab:grid}: with ten components rather than three its
$\fstar$ median is the higher of the two, even though both objectives carry the
same total variance. $\fstar$ sums a per-component maximum
over the active set (\eqref{app:so:a1:eq:fstar}) and each component's
amplitude scales with $\sqrt{s_i}$ by~\eqref{app:so:a1:eq:rescale}, so
spreading one unit of variance over more coordinates raises the sum.

\SOlevelC{\SOcode{aligned10\_sin}}\label{app:so:a2:aligned10-sin}

\textbf{Parameters.} Identical to \SOcode{aligned10} in every field but one:
same $|S|$, same shares, same lengthscale, no pairs, $\gamma=0$, not monotone,
same unit grid --- and the generator is \SOcode{sinusoid} rather than
\SOcode{matern}. The identity goes further than the specification. The
selection stream does not see the generator, so $S$ is the same set at the same
seed, and \SOcode{build} spawns one child stream per active rank the same way
for both, so each component is drawn from the same stream as its
\SOcode{aligned10} counterpart. Only the draw function differs. Its labels
record \SOcode{generator} as \SOcode{sinusoid}, which is how a downstream
analysis tells the two apart.

\textbf{What it isolates.} It is the brief's ``one family generated from
rescaled sinusoids rather than a GP'', one of the two out-of-family checks the
brief names. In every other variant the components are Mat\'ern draws and the
models the study fits are Mat\'ern too, so the additive cells are in-family by
construction and a prior could look good for the wrong reason; this variant
breaks that coincidence while holding everything else fixed. Since
$S$, the shares, the lengthscale and the random stream are all shared with
\SOcode{aligned10}, any difference in a model's behaviour between the two is
attributable to the generator alone.

\textbf{Measured on the archived grid.} Its row of
Table~\ref{app:so:a2:tab:grid}. One caveat belongs with it. The sinusoid
frequency $\omega(\ell)$ of~\eqref{app:so:a1:eq:omega} is chosen to match the
Mat\'ern slope energy per unit variance, but that identity needs at least one
full cycle on the unit interval, $\omega(\ell)\ge1$, and the $\omega$ column of
Table~\ref{app:so:a1:tab:tab_v_ell} puts this variant's lengthscale below that
crossing. For the study's only sinusoid variant the frequency match is
therefore approximate, not exact (\ref{app:so:a1:generators}; design decision
R11,~\ref{app:so:a4}), and its slope shares are read off the recorded label
$g$, a realized quantity, never inferred from~\eqref{app:so:a1:eq:omega}. Its
variance share is unaffected: the rescale of~\eqref{app:so:a1:eq:rescale} is
applied to whatever raw values the generator returns, so $\Varhat f_i=s_i$
holds exactly for it too, by~\eqref{app:so:a1:eq:exactshare}.

\SOlevelC{\SOcode{decoupled}}\label{app:so:a2:decoupled}

\textbf{Parameters.} Eight active coordinates, two to a cell of a
strength-by-roughness grid: two variance shares crossed with two lengthscales,
the four cells Table~\ref{app:so:a2:tab:decoupled} names weak-rough,
weak-smooth, strong-rough and strong-smooth. The four strong and the four weak
shares together sum to one, so $\Varnu f=1$ with no interaction budget to
subtract; there are no pairs, $\gamma=0$, the generator is Mat\'ern, there is no
monotone rejection, and the grid is the unit-interval one.
Table~\ref{app:so:a2:tab:families} gives the shares and lengthscales per rank
rather than as one broadcast value, because a cell needs a specific share paired
with a specific lengthscale and not merely the right multiset --- the
constructor rejects a shares tuple whose length does not match $|S|$ for exactly
that reason. Since shares and lengthscales are assigned by rank, cell membership
is fixed while the coordinates filling the cells are redrawn at every seed.

\textbf{What it isolates.} The design brief's ``the two senses of relevance
disagree''. Variance ranks the coordinates strong above weak and nothing else;
theoretical slope share, by~\eqref{app:so:a1:eq:slopeshare}, weights each share
by $1/(\ell_i^2\vell{\ell_i})$, and at the rough lengthscale that factor is
large enough to lift the weak coordinates over the strong-smooth ones. The
resulting order --- strong-rough, then weak-rough, then strong-smooth, then
weak-smooth --- is the one Table~\ref{app:so:a2:tab:decoupled} prints, and it
crosses the share order: a weak-rough coordinate outranks a strong-smooth one on
slope while the strong-smooth one outranks it on variance. That crossing is the
whole point of the variant. A model that reads relevance off inverse
lengthscales and one that reads it off signal variance cannot both be right
here, so their recovered rankings must separate, and this is the variant where
the thesis expects the two prior families to order the coordinates differently.

\textbf{Measured on the archived grid.} Table~\ref{app:so:a2:tab:decoupled}
gives the theoretical and realized slope shares cell by cell; its row of
Table~\ref{app:so:a2:tab:grid} gives $\fstar$ and the invariants. The realized
column follows the theoretical order across all four cells, which is the check
the design asks of it: the disagreement it manufactures survives the draw.

\begin{table}[htbp]
  \centering
  \caption[The four cells of \SOcode{decoupled}]{
  The four cells of \SOcode{decoupled}: each cell's share and lengthscale, its
  unnormalized theoretical slope energy $s\cdot5/(3\ell^2\vell{\ell})$, the
  theoretical slope share that energy normalizes to over all eight active
  coordinates, and the realized slope share label $g$ averaged over the seeds.
  Mixed check class. Every column but the last is closed form --- a
  deterministic function of the specification and of the
  quadrature~\eqref{app:so:a1:eq:vell} --- and no draw enters it. The last
  column is draw-dependent: measured on
  the archived grid, seeds $0$--$9$, $D=\SOnumStudyD$, this machine's BLAS\@. The
  cell list is derived from the specification's own distinct $(s,\ell)$ pairs,
  so a cell added to the family cannot be silently dropped from the
  table.}\label{app:so:a2:tab:decoupled}
  \small\setlength{\tabcolsep}{3pt}
  \input{\SOroot/tables/tab_decoupled}
\end{table}

\SOlevelC{\SOcode{anti\_aligned}}\label{app:so:a2:anti-aligned}

\textbf{Parameters.} Six active coordinates with shares given per rank in
decreasing order and lengthscales decreasing alongside them --- the brief's
``weak ones rough'', taken to its limit, with one distinct lengthscale per
coordinate. No pairs, $\gamma=0$, Mat\'ern, not monotone, unit grid;
Table~\ref{app:so:a2:tab:families} and Table~\ref{app:so:a2:tab:anti-aligned}
give the values.

This is the one variant whose parameters the design brief's own family table no
longer matches, and here the code is authoritative: the shipped lengthscales are
those the tables print, and the brief's row predates the decision that replaced
them. The reason for the replacement is visible in what the two sets achieve.
With the brief's lengthscales the theoretical slope shares
of~\eqref{app:so:a1:eq:slopeshare} come out neither monotone in rank nor
reversed --- their Spearman correlation against the variance shares is far from
$-1$, and the profile is closer to flat than to reversed --- so the family would
not have delivered the reversal it exists to deliver. The shipped lengthscales
are chosen so that it does, exactly.

\textbf{What it isolates.} The brief's ``sharpest version: share and slope
rankings reversed''. Where \SOcode{decoupled} inverts a single pair of cells ---
weak-rough over strong-smooth --- this variant makes the two senses
anti-correlated over the whole active set: the coordinate with the largest
variance share has the
smallest theoretical slope share and the coordinate with the smallest share has
the largest, with Spearman exactly $-1$ between the two rankings
(Table~\ref{app:so:a2:tab:anti-aligned}). A prior that infers relevance from
inverse lengthscales and one that infers it from signal variance will therefore
rank the coordinates in opposite orders, and a BO loop guided by the wrong one
should spend its budget on the coordinates that matter least. It is the variant
that most sharply separates the thesis's two hypotheses, and the one where the
two are most likely to differ in regret.

\textbf{Measured on the archived grid.}
Table~\ref{app:so:a2:tab:anti-aligned} carries three summary rows, and they are
three different statistics; the entry is worth reading carefully, because they
do not all say the same thing.

The first, the Spearman of the theoretical slope shares against the variance
shares, is a closed-form property of the specification --- no draw enters it ---
and it is exactly $-1$: the design achieves the reversal by construction.

The second is the Spearman of the \emph{seed-mean} realized $g$ against the
shares, the vector of per-rank means over the archived grid's seeds correlated
against $s$. On this grid it is also exactly reversed. This is the statistic the
family's acceptance criterion is stated on.

The third is the Spearman over \emph{all} individual (seed, rank) realized
pairs, pooled into one sample. It is much weaker in magnitude than either of the
others, and that is not a defect in the archive: it is the honest measure of how
well a \emph{single} draw reproduces the design. Realized slope energy is a
draw-level quantity with substantial variance, and this family puts exactly one
coordinate at each lengthscale, so within one seed there is nothing to average
over --- a single unlucky component can invert a neighbouring pair of ranks. The
per-seed realized ranking is therefore unreliable, and no analysis in the thesis
should treat a single seed's realized $g$ order as the family's ground truth;
the reversal is a property of the design and of the seed-averaged behaviour, not
of every draw. This is precisely why the family's acceptance criterion is
stated on the theoretical shares plus a pooled realized check rather than
per seed (design decision R13,~\ref{app:so:a4}). Its row of
Table~\ref{app:so:a2:tab:grid} gives $\fstar$ and the invariants.

\begin{table}[htbp]
  \centering
  \caption[\SOcode{anti\_aligned}: theoretical and realized slope shares]{
  \SOcode{anti\_aligned} rank by rank: the variance share, the lengthscale,
  $\vell{\ell_i}$, the slope factor $5/(3\ell_i^2\vell{\ell_i})$
  of~\eqref{app:so:a1:eq:slopefactor}, the theoretical slope share
  of~\eqref{app:so:a1:eq:slopeshare}, and the realized slope share label $g$
  averaged over the seeds. Mixed check class. Every column but the last is
  closed form, a deterministic function of the specification and of the
  quadrature~\eqref{app:so:a1:eq:vell}, and so is the first summary row. The
  mean realized column and the two remaining summary rows are draw-dependent:
  measured on the archived grid, seeds $0$--$9$, $D=\SOnumStudyD$, this
  machine's BLAS\@. The last two rows are deliberately \emph{different}
  statistics --- the seed-mean vector correlated against $s$, and every
  individual (seed, rank) pair pooled into one sample --- and the second is not
  recoverable from the mean realized column
  alone.}\label{app:so:a2:tab:anti-aligned}
  \small\setlength{\tabcolsep}{3pt}
  \input{\SOroot/tables/tab_anti_aligned}
\end{table}

\SOlevelC{\SOcode{interaction\_g0.00}}\label{app:so:a2:interaction-g0-00}

\textbf{Parameters.} Five active coordinates at one shared lengthscale, plus
two disjoint pairs drawn from inside $S$; Mat\'ern, not monotone, unit grid.
The interaction budget is $\gamma=0$, so each pair carries
$s_{ij}=\gamma/n_{\text{pairs}}=0$ and the main shares are the equal split
$(1-\gamma)/|S|$ of the whole budget. This is the first member of the sweep,
and the four that follow differ from it only in $\gamma$.

It is important that this is \emph{not} simply a five-coordinate aligned
variant. \SOcode{select\_pairs} runs whenever the specification asks for pairs,
irrespective of $\gamma$, so at every seed \SOcode{labels.pairs} records two
pairs; two \SOcode{Interaction} objects are built, each with coefficient
$c_{ij}=\sqrt{s_{ij}}=0$ and each with its own two unit-variance factor draws;
and $\fstar$ is assembled by the two-dimensional block search over those pairs
rather than by the per-component maximum, since the assembly treats a
coordinate in any pair as part of a block. The variant is the sweep's control,
sharing the sweep's code path and its random stream exactly, and its recorded
labels look like an interacting objective whose interactions happen to have
zero weight.

\textbf{What it isolates.} The sweep as a whole is the design brief's
``additivity failing''; this member is its additive end, the reference the
other four are read against. Because it runs the same machinery at zero
strength, any difference between it and a member with $\gamma>0$ is
attributable to the interaction budget alone --- not to a change of code path,
of active set, of pair choice, or of random stream. It is also the variant that
tells the study whether a purely additive model pays any price merely for the
possibility of interaction.

\textbf{Measured on the archived grid.} Its row of
Table~\ref{app:so:a2:tab:grid}. It is the only member of the sweep whose QMC
$\Varnu$ minimum and maximum agree to the printed precision: with $\gamma=0$
every $c_{ij}$ is zero, so the objective is exactly additive and the pair
blocks the other four rows sample over contribute nothing.

\SOlevelC{\SOcode{interaction\_g0.10}}\label{app:so:a2:interaction-g0-10}

\textbf{Parameters.} The same shape as \SOcode{interaction\_g0.00} --- five
active coordinates, one lengthscale, two disjoint pairs, Mat\'ern, not
monotone, unit grid --- with the interaction budget $\gamma$ raised to the
value Table~\ref{app:so:a2:tab:families} prints, split equally between the two
pairs as $s_{ij}=\gamma/n_{\text{pairs}}$ with $c_{ij}=\sqrt{s_{ij}}$
(\eqref{app:so:a1:eq:hshare}). The main shares fall to $(1-\gamma)/|S|$ so that
$\sum_i s_i+\gamma=1$ and $\Varnu f$ is still one
(\eqref{app:so:a1:eq:budget}). At a given seed the active set, the two pairs
and every raw draw are the same as in every other member of the sweep; only the
shares and the pair coefficients move.

\textbf{What it isolates.} ``Additivity failing'', at the sweep's smallest
nonzero dose. A tenth of the variance sitting outside the additive
decomposition is the regime where an additive model is still nearly right, so
this is where the thesis expects a well-calibrated sparsity prior to lose
little: the mismatch should show up as a small, systematic gap rather than a
failure of support recovery.

\textbf{Measured on the archived grid.} Its row of
Table~\ref{app:so:a2:tab:grid}.

\SOlevelC{\SOcode{interaction\_g0.25}}\label{app:so:a2:interaction-g0-25}

\textbf{Parameters.} As \SOcode{interaction\_g0.10} with $\gamma$ at the next
value in Table~\ref{app:so:a2:tab:families}: five active coordinates, one
lengthscale, two disjoint pairs sharing $\gamma$ equally, main shares
$(1-\gamma)/|S|$. Active set, pairs and raw draws are again unchanged from the
rest of the sweep at the same seed.

\textbf{What it isolates.} ``Additivity failing'' at a quarter of the variance.
This is the middle of the sweep, where the interaction is large enough that an
additive surrogate must misfit somewhere but small enough that the main effects
still dominate the ranking; it is the member most likely to separate a model
that degrades gracefully under misspecification from one that does not.

\textbf{Measured on the archived grid.} Its row of
Table~\ref{app:so:a2:tab:grid}.

\SOlevelC{\SOcode{interaction\_g0.50}}\label{app:so:a2:interaction-g0-50}

\textbf{Parameters.} As the rest of the sweep, with $\gamma$ at the value
Table~\ref{app:so:a2:tab:families} prints: half the variance budget in the two
pairs and half spread equally over the five main effects, each of which now
carries a smaller share than its pair does.

\textbf{What it isolates.} ``Additivity failing'' at the point where the two
halves of the decomposition are balanced. Beyond here the additive part of the
objective is no longer the larger part, so a model that can only represent main
effects is fitting a minority of the signal --- and, because the pairs sit
inside $S$, it is fitting it on coordinates that are genuinely relevant, which
is the failure mode the sweep is built to grade.

\textbf{Measured on the archived grid.} Its row of
Table~\ref{app:so:a2:tab:grid}.

\SOlevelC{\SOcode{interaction\_g0.75}}\label{app:so:a2:interaction-g0-75}

\textbf{Parameters.} The sweep's largest interaction budget, again split
equally between the same two pairs, with the five main shares reduced to
$(1-\gamma)/|S|$ accordingly; every other field is as
\SOcode{interaction\_g0.00}.

\textbf{What it isolates.} ``Additivity failing'' taken to the point where most
of the variance is in the pairs. An additive model is now badly misspecified by
construction, and the question the variant asks is not whether it fits but
whether the sparsity prior still recovers the right \emph{support} --- the
first-order shares $s_i$ remain equal and nonzero, so the active set is
unchanged even though most of the function is no longer additive.

\textbf{Measured on the archived grid.} Its row of
Table~\ref{app:so:a2:tab:grid}, and, read together with the four rows above it,
the sweep's trend: the median $\fstar$ rises monotonically with $\gamma$ and
the spread between the smallest and largest $\fstar$ across seeds widens
monotonically with it, this variant having the widest of the five. Both follow
from the product form of~\eqref{app:so:a1:eq:hdef}: a pair block's maximum is
the largest value of $c_{ij}u_i u_j$ over the square, so moving budget into the
pairs both raises the attainable maximum and makes it depend on the shape of
two draws at once rather than one.

\SOlevelC{\SOcode{rotated\_t0}}\label{app:so:a2:rotated-t0}

\textbf{Parameters.} Ten active coordinates with equal shares at one
lengthscale, arranged in five disjoint pairs, each pair rotated by $\theta$ in
its own plane (\ref{app:so:a1:rotation}); Mat\'ern, not monotone. There is no
explicit interaction: $\gamma$ is exactly zero and no \SOcode{Interaction}
object is built at all, because with $\theta$ set the pairs are what the
rotation acts on and the shares alone already sum to one. Every component sits
on the extended grid, $\knots=\SOnumKnotsExt$ rather than $\SOnumKnotsUnit$
(Table~\ref{app:so:a2:tab:families}), since a rotated coordinate's argument
runs outside the unit interval (\eqref{app:so:a1:eq:next}).

That last point has a consequence worth stating plainly, because it is easy to
assume otherwise: at $\theta=0$ the rotation is the identity and the objective
\emph{is} exactly additive and axis-aligned, but it is not the same objective
as \SOcode{aligned10} at the same seed. \SOcode{build} chooses the draw domain
and the grid size per coordinate according to whether that coordinate sits in a
rotation pair, so every component here is drawn on the extended grid instead of
the unit one --- a different eigen factor and a different number of variates ---
and the delivered functions differ. The active set is the same; the draws are
not.

\textbf{What it isolates.} The rotated family is the design brief's ``axis
alignment failing''; this member is its aligned end, the control that separates
the effect of rotating from the effect of being built for rotation. Comparing it
with \SOcode{rotated\_t15} and \SOcode{rotated\_t45} isolates the angle, since
those three differ in $\theta$ and in nothing else; comparing it with
\SOcode{aligned10} does not, because that comparison confounds the angle with
the change of draw domain described above.

\textbf{Measured on the archived grid.} Its row of
Table~\ref{app:so:a2:tab:grid} and the $\theta=0$ row of
Table~\ref{app:so:a1:tab:tab_rotation_angle}, where $\kappa$ is exactly one at
every seed and $\gaxis$ is zero up to floating-point epsilon --- never a
genuinely negative quantity; that table's own header comment records the
per-seed residuals. Both repairs the rotation needs, the recentering $\mu$ and
the renormalization $\kappa$, are therefore inactive here.

\SOlevelC{\SOcode{rotated\_t15}}\label{app:so:a2:rotated-t15}

\textbf{Parameters.} As \SOcode{rotated\_t0} in every field but $\theta$, which
Table~\ref{app:so:a2:tab:families} gives: ten coordinates, equal shares, one
lengthscale, five disjoint rotated pairs, $\gamma$ exactly zero, extended grid.

\textbf{What it isolates.} ``Axis alignment failing'' at a small angle. The
function is still additive in a rotated basis and still exactly the same
function of that basis; what fails is only the assumption that the relevant
directions are coordinate axes. This is the variant where an axis-aligned
sparsity prior is mildly wrong, and the thesis's hypothesis is that the two
prior families differ in how gracefully they absorb that: a small rotation
leaks a little variance into pairwise structure that no per-coordinate
lengthscale can represent.

\textbf{Measured on the archived grid.} Its row of
Table~\ref{app:so:a2:tab:grid} and the $\theta=15^\circ$ row of
Table~\ref{app:so:a1:tab:tab_rotation_angle}. Two things there deserve care.
First, $\kappa$ straddles one at this angle: the table's count of seeds with
$\kappa>1$ is nonzero here, so the sign of $\kappa-1$ is seed-dependent at
$15^\circ$ --- as it also is at $60^\circ$, and as it is not at $30^\circ$ or
$45^\circ$, where every seed has $\kappa<1$. Second, the angle does not index
the misspecification; see the entry for \SOcode{rotated\_t45}.

\SOlevelC{\SOcode{rotated\_t45}}\label{app:so:a2:rotated-t45}

\textbf{Parameters.} As \SOcode{rotated\_t0} and \SOcode{rotated\_t15} but at
the largest of the study's three angles: ten coordinates, equal shares, one
lengthscale, five disjoint rotated pairs, $\gamma$ exactly zero, extended grid.

\textbf{What it isolates.} ``Axis alignment failing'' at the diagonal, where a
pair's two rotated directions are equal mixtures of its two coordinates and the
axis-aligned decomposition has least to work with. Of the six angles
Table~\ref{app:so:a1:tab:tab_rotation_angle} reports, this one has the largest
median $\gaxis$. It is the rotated family's worst case, and the natural partner
to \SOcode{anti\_aligned}: both ask what happens when the structure the prior
assumes is not the structure the objective has.

\textbf{Measured on the archived grid.} Its row of
Table~\ref{app:so:a2:tab:grid} and the $\theta=45^\circ$ row of
Table~\ref{app:so:a1:tab:tab_rotation_angle}. Its QMC $\Varnu$ minimum is the
lowest in the grid summary; the final row of that table repeats this variant at
one seed with a sixteen-fold larger Sobol' sample, and both moment columns
tighten, so that residual reads as sampling error rather than a defect in the
objective.

The analysis rule for the whole rotated family (\ref{app:so:a1:rotation},
ruling R28) is best stated here, since this is the angle at which it bites.
$\gaxis$ vanishes at every multiple of $90^\circ$ and is largest near the
diagonal, so it is not monotone in $\theta$; and the $\gaxis$ ranges at the
$30^\circ$, $45^\circ$ and $60^\circ$ rows overlap one another, so among those
rows the nominal angle does not determine $\gaxis$ --- restricted to those three
rows, the spread across seeds within one angle exceeds the gap between the
medians of adjacent angles. The study therefore conditions its analysis on the
recorded \SOcode{gamma\_axis} label of each objective, never on $\theta$. The
nominal angle names the variant; it does not measure it.

\SOlevelC{\SOcode{dense\_weak}}\label{app:so:a2:dense-weak}

\textbf{Parameters.} Every one of the $D$ coordinates is active, each with the
same share $1/D$, all at one large lengthscale; no pairs, $\gamma=0$, Mat\'ern
generator, and --- alone among the fourteen --- the monotone rejection wrapper
is on, so each component is redrawn until its grid values are monotone
(\ref{app:so:a1:generators}). The direction is drawn per coordinate, so a
component is as likely to be decreasing as increasing. A component whose values
rise or fall monotonically across the grid takes its largest value at an end of
the interval, and on the archived grid this holds without exception: at every
seed the maximizer $x^\star$ the labels record is a corner of the cube, every
coordinate at exactly $0$ or $1$, with both directions represented. The explicit
direction draw is in fact redundant in law, as shown
in~\ref{app:so:a1:generators} --- the rejection already accepts either direction
--- so the corner structure follows from the monotonicity and not from the coin.

This is the one variant whose row in Table~\ref{app:so:a2:tab:families} is not
the study's own configuration: $|S|$ and the share both track $D$, and that
table is built at $D=\SOnumTableD$ rather than at the study's
$D=\SOnumStudyD$. The shape is the same at either dimension --- every
coordinate active with share $1/D$ --- but the numbers in that one row are not
the study's, and the next paragraph turns on the difference.

\textbf{What it isolates.} The design brief's ``sparsity failing
(Lasso-DNA-like); all shares below $\activeeps$ by design''. It is the arm where
the sparse prior's central assumption is false: there is no small relevant
subset to find, only many weak coordinates. The thesis's hypothesis is that a
prior which shrinks aggressively toward sparsity should pay here, while one
that can express many weak relevances should not --- and, because every
coordinate is monotone at a large lengthscale, the objective remains easy for a
method that simply follows the gradient of the mean, so a failure on this
variant is a failure of the prior and not of the problem's difficulty.

The ``below $\activeeps$'' half of the brief's description is worth stating
exactly, because it is a property of $D$ and not of the family. The share is
$1/D$, so $1/D<\activeeps$ exactly when $D>1/\activeeps$: at the study's
$D=\SOnumStudyD$ every share is below the threshold, which is the design point,
while at the $D=\SOnumTableD$ Table~\ref{app:so:a2:tab:families} is built at,
no share is. Either way \SOcode{Labels.active} is all-true, at every $D$: the
label is set membership, $s_i>0$, and $\activeeps$ is recorded separately in
\SOcode{Labels.active\_eps} for scoring rather than enforced. The variant is
built precisely so that the brief's scoring rule and the objective's own notion
of an active coordinate disagree at the study's dimension.

\textbf{Measured on the archived grid.} Its row of
Table~\ref{app:so:a2:tab:grid}. It has by a wide margin the largest $\fstar$ in
the table and the narrowest spread of $\fstar$ across seeds --- both direct
consequences of its shape: $\fstar$ sums a maximum over all $D$ coordinates
rather than a handful, and averaging that many independent per-component maxima
leaves little seed-to-seed variation.
